% =====================================================================
% Abbildung: Traceability-Kette am Beispiel REQ-002
% Passend zu Kapitel 6, Abschnitt "SysML-Beziehungen"
%
% Einbindung im Buch, z. B. in chapters/06_anforderungen.tex:
%
%   \begin{figure}[H]
%     \centering
%     % =====================================================================
% Abbildung: Traceability-Kette am Beispiel REQ-002
% Passend zu Kapitel 6, Abschnitt "SysML-Beziehungen"
%
% Einbindung im Buch, z. B. in chapters/06_anforderungen.tex:
%
%   \begin{figure}[H]
%     \centering
%     % =====================================================================
% Abbildung: Traceability-Kette am Beispiel REQ-002
% Passend zu Kapitel 6, Abschnitt "SysML-Beziehungen"
%
% Einbindung im Buch, z. B. in chapters/06_anforderungen.tex:
%
%   \begin{figure}[H]
%     \centering
%     % =====================================================================
% Abbildung: Traceability-Kette am Beispiel REQ-002
% Passend zu Kapitel 6, Abschnitt "SysML-Beziehungen"
%
% Einbindung im Buch, z. B. in chapters/06_anforderungen.tex:
%
%   \begin{figure}[H]
%     \centering
%     \input{figures/fig_kap06_traceability_req002}
%     \caption{Traceability-Kette von REQ-002 ueber \texttt{satisfy} und
%       \texttt{verify} zum \texttt{PerceptionSystem} und per \texttt{trace}
%       weiter zum ROS2-Node.}
%     \label{fig:kap06-traceability-req002}
%   \end{figure}
%
% Benoetigte Pakete (im Buch bereits in main.tex geladen): tikz
% Benoetigte TikZ-Bibliotheken: positioning, arrows.meta, shadows, calc
% =====================================================================

\input{figures/fig_kap06_traceability_req002.tikz}

%     \caption{Traceability-Kette von REQ-002 ueber \texttt{satisfy} und
%       \texttt{verify} zum \texttt{PerceptionSystem} und per \texttt{trace}
%       weiter zum ROS2-Node.}
%     \label{fig:kap06-traceability-req002}
%   \end{figure}
%
% Benoetigte Pakete (im Buch bereits in main.tex geladen): tikz
% Benoetigte TikZ-Bibliotheken: positioning, arrows.meta, shadows, calc
% =====================================================================

\begin{tikzpicture}[
    every node/.style={font=\small},
    req/.style={
        rectangle, rounded corners=1.5pt, draw=black!70, thick,
        minimum width=34mm, minimum height=12mm,
        align=center, fill=blue!12, drop shadow
    },
    part/.style={
        rectangle, rounded corners=1.5pt, draw=green!40!black, thick,
        minimum width=34mm, minimum height=12mm,
        align=center, fill=green!15, drop shadow
    },
    test/.style={
        rectangle, rounded corners=1.5pt, draw=black!50, thick,
        minimum width=34mm, minimum height=12mm,
        align=center, fill=black!6, drop shadow
    },
    code/.style={
        rectangle, rounded corners=1.5pt, draw=orange!60!black, thick,
        minimum width=40mm, minimum height=12mm,
        align=center, fill=orange!15, drop shadow
    },
    flow/.style={-{Stealth[length=2.2mm]}, thick, black!60},
    lbl/.style={font=\small, fill=white, inner sep=1pt}
]

% --- Knoten ---
\node[req]  (req002)  at (0, 3)     {\texttt{REQ-002}\\Hindernisse erkennen};
\node[part] (percept) at (7, 3)     {\texttt{PerceptionSystem}};
\node[test] (tc002)   at (0, -1.5)  {\texttt{TC-002}\\Hinderniserkennung};
\node[code] (rosnode) at (7, -1.5)  {\texttt{obstacle\_detector\_node}};

% --- Beziehungen ---
\draw[flow] (req002) -- node[lbl, midway]{\texttt{satisfy}} (percept);
\draw[flow] (tc002)  to[bend left=12] node[lbl, near end]{\texttt{verify}} (percept);
\draw[flow] (percept) -- node[lbl, midway]{\texttt{trace}} (rosnode);

% --- Legende ---
\node[req,  minimum width=6mm, minimum height=4mm, below=15mm of tc002, xshift=-8mm, font=\tiny] (leg1) {};
\node[right=1mm of leg1, font=\tiny, align=left] {Anforderung};
\node[part, minimum width=6mm, minimum height=4mm, below=6mm of leg1, font=\tiny] (leg2) {};
\node[right=1mm of leg2, font=\tiny, align=left] {Systemteil (\texttt{part def})};
\node[test, minimum width=6mm, minimum height=4mm, below=6mm of leg2, font=\tiny] (leg3) {};
\node[right=1mm of leg3, font=\tiny, align=left] {Testfall};
\node[code, minimum width=6mm, minimum height=4mm, below=6mm of leg3, font=\tiny] (leg4) {};
\node[right=1mm of leg4, font=\tiny, align=left] {ROS2-Node};

\end{tikzpicture}


%     \caption{Traceability-Kette von REQ-002 ueber \texttt{satisfy} und
%       \texttt{verify} zum \texttt{PerceptionSystem} und per \texttt{trace}
%       weiter zum ROS2-Node.}
%     \label{fig:kap06-traceability-req002}
%   \end{figure}
%
% Benoetigte Pakete (im Buch bereits in main.tex geladen): tikz
% Benoetigte TikZ-Bibliotheken: positioning, arrows.meta, shadows, calc
% =====================================================================

\begin{tikzpicture}[
    every node/.style={font=\small},
    req/.style={
        rectangle, rounded corners=1.5pt, draw=black!70, thick,
        minimum width=34mm, minimum height=12mm,
        align=center, fill=blue!12, drop shadow
    },
    part/.style={
        rectangle, rounded corners=1.5pt, draw=green!40!black, thick,
        minimum width=34mm, minimum height=12mm,
        align=center, fill=green!15, drop shadow
    },
    test/.style={
        rectangle, rounded corners=1.5pt, draw=black!50, thick,
        minimum width=34mm, minimum height=12mm,
        align=center, fill=black!6, drop shadow
    },
    code/.style={
        rectangle, rounded corners=1.5pt, draw=orange!60!black, thick,
        minimum width=40mm, minimum height=12mm,
        align=center, fill=orange!15, drop shadow
    },
    flow/.style={-{Stealth[length=2.2mm]}, thick, black!60},
    lbl/.style={font=\small, fill=white, inner sep=1pt}
]

% --- Knoten ---
\node[req]  (req002)  at (0, 3)     {\texttt{REQ-002}\\Hindernisse erkennen};
\node[part] (percept) at (7, 3)     {\texttt{PerceptionSystem}};
\node[test] (tc002)   at (0, -1.5)  {\texttt{TC-002}\\Hinderniserkennung};
\node[code] (rosnode) at (7, -1.5)  {\texttt{obstacle\_detector\_node}};

% --- Beziehungen ---
\draw[flow] (req002) -- node[lbl, midway]{\texttt{satisfy}} (percept);
\draw[flow] (tc002)  to[bend left=12] node[lbl, near end]{\texttt{verify}} (percept);
\draw[flow] (percept) -- node[lbl, midway]{\texttt{trace}} (rosnode);

% --- Legende ---
\node[req,  minimum width=6mm, minimum height=4mm, below=15mm of tc002, xshift=-8mm, font=\tiny] (leg1) {};
\node[right=1mm of leg1, font=\tiny, align=left] {Anforderung};
\node[part, minimum width=6mm, minimum height=4mm, below=6mm of leg1, font=\tiny] (leg2) {};
\node[right=1mm of leg2, font=\tiny, align=left] {Systemteil (\texttt{part def})};
\node[test, minimum width=6mm, minimum height=4mm, below=6mm of leg2, font=\tiny] (leg3) {};
\node[right=1mm of leg3, font=\tiny, align=left] {Testfall};
\node[code, minimum width=6mm, minimum height=4mm, below=6mm of leg3, font=\tiny] (leg4) {};
\node[right=1mm of leg4, font=\tiny, align=left] {ROS2-Node};

\end{tikzpicture}


%     \caption{Traceability-Kette von REQ-002 ueber \texttt{satisfy} und
%       \texttt{verify} zum \texttt{PerceptionSystem} und per \texttt{trace}
%       weiter zum ROS2-Node.}
%     \label{fig:kap06-traceability-req002}
%   \end{figure}
%
% Benoetigte Pakete (im Buch bereits in main.tex geladen): tikz
% Benoetigte TikZ-Bibliotheken: positioning, arrows.meta, shadows, calc
% =====================================================================

\begin{tikzpicture}[
    every node/.style={font=\small},
    req/.style={
        rectangle, rounded corners=1.5pt, draw=black!70, thick,
        minimum width=34mm, minimum height=12mm,
        align=center, fill=blue!12, drop shadow
    },
    part/.style={
        rectangle, rounded corners=1.5pt, draw=green!40!black, thick,
        minimum width=34mm, minimum height=12mm,
        align=center, fill=green!15, drop shadow
    },
    test/.style={
        rectangle, rounded corners=1.5pt, draw=black!50, thick,
        minimum width=34mm, minimum height=12mm,
        align=center, fill=black!6, drop shadow
    },
    code/.style={
        rectangle, rounded corners=1.5pt, draw=orange!60!black, thick,
        minimum width=40mm, minimum height=12mm,
        align=center, fill=orange!15, drop shadow
    },
    flow/.style={-{Stealth[length=2.2mm]}, thick, black!60},
    lbl/.style={font=\small, fill=white, inner sep=1pt}
]

% --- Knoten ---
\node[req]  (req002)  at (0, 3)     {\texttt{REQ-002}\\Hindernisse erkennen};
\node[part] (percept) at (7, 3)     {\texttt{PerceptionSystem}};
\node[test] (tc002)   at (0, -1.5)  {\texttt{TC-002}\\Hinderniserkennung};
\node[code] (rosnode) at (7, -1.5)  {\texttt{obstacle\_detector\_node}};

% --- Beziehungen ---
\draw[flow] (req002) -- node[lbl, midway]{\texttt{satisfy}} (percept);
\draw[flow] (tc002)  to[bend left=12] node[lbl, near end]{\texttt{verify}} (percept);
\draw[flow] (percept) -- node[lbl, midway]{\texttt{trace}} (rosnode);

% --- Legende ---
\node[req,  minimum width=6mm, minimum height=4mm, below=15mm of tc002, xshift=-8mm, font=\tiny] (leg1) {};
\node[right=1mm of leg1, font=\tiny, align=left] {Anforderung};
\node[part, minimum width=6mm, minimum height=4mm, below=6mm of leg1, font=\tiny] (leg2) {};
\node[right=1mm of leg2, font=\tiny, align=left] {Systemteil (\texttt{part def})};
\node[test, minimum width=6mm, minimum height=4mm, below=6mm of leg2, font=\tiny] (leg3) {};
\node[right=1mm of leg3, font=\tiny, align=left] {Testfall};
\node[code, minimum width=6mm, minimum height=4mm, below=6mm of leg3, font=\tiny] (leg4) {};
\node[right=1mm of leg4, font=\tiny, align=left] {ROS2-Node};

\end{tikzpicture}

